%! Function Fundamentals: Notation, Domain, and Range — Unit 2, Lesson 2.1 — Student Packet Cover Sheet
\documentclass[10pt]{article}
\usepackage{atda-article}
\usepackage{atda-boxes}
\usepackage{ltablex}
\keepXColumns

\begin{document}

% Full-bleed royal banner
\begin{tikzpicture}[remember picture, overlay]
  \fill[royal] (current page.north west)
    rectangle ([yshift=-0.9in]current page.north east);
\end{tikzpicture}
\vspace{-0.8in}

\begin{center}
  {\color{white}\textbf{\LARGE Algebra, Trigonometry, and Data Analysis}}\\[4pt]
  {\color{mistmid}\large\textbf{Unit 2: Functions, Transformations, and Graph Analysis}}\\[6pt]
  {\color{white}\textbf{\large Lesson 2.1 \quad Function Fundamentals: Notation, Domain, and Range}}
\end{center}

\vspace{0.22in}
\namedateperiod
\vspace{0.18in}

\begin{learningtargetbox}
\small
\begin{enumerate}[leftmargin=1.5em, itemsep=5pt]
  \item \ldots{} read a graph in \textbf{both directions} --- given an input, find the output
        $f(x)$; given an output, find the input that produced it --- and write either read as a
        point $(x, f(x))$.
  \item \ldots{} state the \textbf{domain} and \textbf{range} of a function from its graph, with
        units, in inequality notation, set notation, and interval notation.
  \item \ldots{} explain what cuts a domain short: the situation being modeled, where the graph
        stops, or a value the algebra will not allow.
  \item \ldots{} turn a point on a graph into a sentence about the situation, and justify why one
        output can come from two different inputs but one input can never give two outputs.
\end{enumerate}
\end{learningtargetbox}

\vspace{0.14in}

\begin{tocbox}
\small
\renewcommand{\arraystretch}{1.5}
\begin{tabularx}{\linewidth}{@{} c l X r @{}}
\rowcolor{mist}
\textbf{\#} & \textbf{Component} & \textbf{Description} & \textbf{Score} \\
\midrule
1 & Warm-Up        & Function notation both ways, an excluded value, and writing an interval
                                                                    & \blank{1.2cm} \\
2 & Guided Notes   & The parking garage: two questions, one graph --- then domain, range, and the
                     three ways to write them                       & \blank{1.2cm} \\
3 & Group Activity & A burning candle, a skate bowl, and a bus that only takes whole students
                                                                    & \blank{1.2cm} \\
4 & Exit Ticket    & Reading a draining pool in both directions, with domain and range in context
                                                                    & \blank{1.2cm} \\
5 & Homework       & Notation practice, a highway arch, and a look ahead & \blank{1.2cm} \\
\midrule
  & \multicolumn{2}{r}{\textbf{Total}} & \blank{1.2cm} \\
\end{tabularx}
\end{tocbox}

\vspace{0.14in}

\begin{remindbox}
\small
This lesson covers \texttt{AFDA.AF.2a} (determine the domain and range of a function from its
graph, \emph{including those limited by contexts}) and \texttt{AFDA.AF.2e} (the connection between
a point on the graph and the value of the function). Two sentences carry the whole lesson:
\textbf{``$f(3)=10$'' and ``the point $(3,10)$ is on the graph'' say exactly the same thing}, and
\textbf{a model is only in force where the situation is}. Bring Lesson 1.7 with you --- an excluded
value is about to become a hole in a domain.
\end{remindbox}

\end{document}
